\documentclass[a4paper,11pt,fleqn]{scrartcl}%fleqn pour aligner les equations à gauche
\usepackage{../../Styles/Exercice}


\newcommand{\chnum}{Chapitre 14}
\newcommand{\chtitle}{Modélisation de l'écoulement d'un fluide}
\newcommand{\dtype}{Exercices}
  
\begin{document}
\section*{Corrigé}

\section{Comprendre le déplacement vertical des poissons}

\begin{enumerate}
\item Lorsque le poisson gonfle ou dégonfle sa vessie natatoire, la grandeur physique qui varie est le \textbf{volume} du poisson. Sa masse reste pratiquement constante mais son volume change.

\item Bilan des forces exercées sur un poisson immobile dans l'eau :
\begin{itemize}
\item son \textbf{poids} $\vec{P}=m\vec{g}$ dirigé vers le bas ;
\item la \textbf{poussée d'Archimède} $\vec{F_A}$ dirigée vers le haut.
\end{itemize}

À l'équilibre (poisson immobile) :

\[
F_A = P
\]

Si le poisson gonfle sa vessie natatoire, son volume augmente. La poussée d'Archimède

\[
F_A = \rho_{eau} g V
\]

augmente donc également. On obtient alors :

\[
F_A > P
\]

La résultante des forces est dirigée vers le haut : le poisson \textbf{remonte}.
\end{enumerate}

\section{Calculer une norme}

\begin{enumerate}

\item La poussée d'Archimède vaut :

\[
F_A = \rho_{eau} g V
\]

Avec :

\[
\rho_{eau}= \SI{1000}{kg.m^{-3}}, \quad g=\SI{9,81}{m.s^{-2}}, \quad V=\SI{2,0}{m^3}
\]

\[
F_A = 1000 \times 9,81 \times 2,0
\]

\[
F_A = \SI{1,96e4}{N}
\]

La poussée d'Archimède exercée sur la barque vaut donc environ :

\[
F_A \approx \SI{2,0e4}{N}
\]

\item La poussée d'Archimède provient du fait que la pression du fluide augmente avec la profondeur.  
La pression exercée sur la partie inférieure de l'objet immergé est donc plus grande que celle exercée sur sa partie supérieure.  
Cette différence de pression crée une force résultante verticale dirigée vers le haut.
\end{enumerate}

\section{Calculer un débit volumique}

\begin{enumerate}

\item Le débit volumique est défini par :

\[
Q = \frac{V}{\Delta t}
\]

\[
V=\SI{1,5}{L}= \SI{1,5e-3}{m^3}
\]

\[
\Delta t = \SI{60}{s}
\]

\[
Q=\frac{1,5\times10^{-3}}{60}
\]

\[
Q=\SI{2,5e-5}{m^3.s^{-1}}
\]

\item La relation entre débit et vitesse est :

\[
Q = S v
\]

avec

\[
S = \SI{1,0}{cm^2} = \SI{1,0e-4}{m^2}
\]

\[
v = \frac{Q}{S}
\]

\[
v = \frac{2,5\times10^{-5}}{1,0\times10^{-4}}
\]

\[
v = \SI{0,25}{m.s^{-1}}
\]

La vitesse d'écoulement de l'eau est donc :

\[
v = \SI{0,25}{m.s^{-1}}
\]

\end{enumerate}

\section{Calculer une dépression}

\begin{enumerate}

\item Équation de continuité :

\[
Q = S_1 v_1 = S_2 v_2
\]

\[
v_2 = \frac{S_1}{S_2} v_1
\]

\[
v_2 = \frac{2,0}{1,0} \times 5,0
\]

\[
v_2 = \SI{10}{m.s^{-1}}
\]

\item Équation de Bernoulli pour un écoulement horizontal :

\[
P_1 + \frac12 \rho v_1^2 = P_2 + \frac12 \rho v_2^2
\]

On obtient :

\[
P_1 - P_2 = \frac12 \rho (v_2^2 - v_1^2)
\]

\[
P_1 - P_2 = \frac12 \times 1,2 \times (10^2 - 5^2)
\]

\[
P_1 - P_2 = 0,6 \times 75
\]

\[
P_1 - P_2 = \SI{45}{Pa}
\]

La dépression vaut donc :

\[
\Delta P = \SI{45}{Pa}
\]

\end{enumerate}

\section{Calculer une vitesse d'éjection de vidange}

On applique le théorème de Bernoulli entre la surface libre et l'orifice.

La pression est atmosphérique aux deux points et la vitesse de la surface est négligeable.

On obtient la \textbf{loi de Torricelli} :

\[
v = \sqrt{2 g h_0}
\]

Avec :

\[
h_0=\SI{1,0}{m}
\]

\[
v = \sqrt{2\times 9,81 \times 1,0}
\]

\[
v = \sqrt{19,62}
\]

\[
v \approx \SI{4,4}{m.s^{-1}}
\]

\section{Siphon}

\begin{enumerate}

\item Comme $s << \Sigma$, le débit étant conservé :

\[
Q = s v_C = \Sigma v_D
\]

Donc :

\[
v_D << v_C
\]

La vitesse au point $D$ est donc \textbf{négligeable} devant celle au point $C$.

\item Bernoulli entre $D$ et $C$ :

\[
P_a + \rho g z_D = P_a + \frac12 \rho v_C^2 + \rho g z_C
\]

\[
\frac12 \rho v_C^2 = \rho g (z_D - z_C)
\]

\[
v_C = \sqrt{2 g (z_D - z_C)}
\]

\[
z_D - z_C = 0,60 - (-0,10) = 0,70
\]

\[
v_C = \sqrt{2\times 9,81\times0,70}
\]

\[
v_C \approx \SI{3,7}{m.s^{-1}}
\]

Condition d'écoulement :

\[
z_D > z_C
\]

\item Pression au point $A$

Bernoulli entre $D$ et $A$ :

\[
P_a + \rho g z_D = P_A + \frac12 \rho v_A^2 + \rho g z_A
\]

avec $v_A = v_C$.

On obtient :

\[
P_A = P_a + \rho g (z_D - z_A) - \frac12 \rho v_C^2
\]

Pression au point $B$ :

\[
P_B = P_a + \rho g (z_D - z_B) - \frac12 \rho v_C^2
\]

\item Pour amorcer le siphon, il faut \textbf{remplir le tube d'eau} afin d'éviter la présence d'air.

La hauteur du point $B$ ne peut pas être quelconque.  
Si $B$ est trop haut, la pression $P_B$ devient trop faible et peut atteindre la pression de vapeur de l'eau : des bulles apparaissent et le siphon cesse de fonctionner.
\end{enumerate}

\section{Tube de Venturi — Corrigé}

\begin{enumerate}

\item \textbf{Différence de pression entre les positions 1 et 2}

\textbf{a) Expression en fonction de la différence de niveaux $h$}

Dans le tube en U, les fluides sont au repos : on applique la relation fondamentale de l'hydrostatique.

La différence de pression entre les points 1 et 2 est liée à la différence de niveaux du mercure :

\[
P_1 - P_2 = (\rho_{\ce{Hg}}-\rho_{eau}) g h
\]

Cette relation tient compte du fait que les deux branches contiennent au-dessus du mercure de l'eau.

\bigskip

\textbf{b) Expression en fonction du débit volumique}

On applique l'équation de Bernoulli entre les sections 1 et 2 du tube de Venturi (écoulement horizontal) :

\[
P_1 + \frac12 \rho_{eau} v_1^2
=
P_2 + \frac12 \rho_{eau} v_2^2
\]

d'où :

\[
P_1 - P_2 = \frac12 \rho_{eau}(v_2^2-v_1^2)
\]

Le débit volumique est constant :

\[
D_V = S_1 v_1 = S_2 v_2
\]

avec

\[
S = \frac{\pi d^2}{4}
\]

Donc

\[
v_1 = \frac{D_V}{S_1}, \qquad v_2 = \frac{D_V}{S_2}
\]

On obtient :

\[
P_1-P_2 =
\frac12 \rho_{eau} D_V^2
\left(
\frac{1}{S_2^2}-\frac{1}{S_1^2}
\right)
\]

\bigskip

\item \textbf{Débit volumique en fonction de $h$}

On identifie les deux expressions de la différence de pression :

\[
(\rho_{\ce{Hg}}-\rho_{eau})gh
=
\frac12 \rho_{eau} D_V^2
\left(
\frac{1}{S_2^2}-\frac{1}{S_1^2}
\right)
\]

On isole le débit volumique :

\[
D_V =
\sqrt{
\frac{2(\rho_{\ce{Hg}}-\rho_{eau})gh}
{\rho_{eau}
\left(\frac{1}{S_2^2}-\frac{1}{S_1^2}\right)}
}
\]

\bigskip

\textbf{Calcul des sections}

\[
d_1=\SI{10}{mm}=1,0\times10^{-2}\,\text{m}
\]

\[
d_2=\SI{5,0}{mm}=5,0\times10^{-3}\,\text{m}
\]

\[
S_1=\frac{\pi d_1^2}{4}
=7,85\times10^{-5}\,\text{m}^2
\]

\[
S_2=\frac{\pi d_2^2}{4}
=1,96\times10^{-5}\,\text{m}^2
\]

La relation finale devient donc :

\[
D_V =
\sqrt{
\frac{2(13550-1000)gh}
{1000
\left(\frac{1}{S_2^2}-\frac{1}{S_1^2}\right)}
}
\]

Cette expression permet de déterminer expérimentalement le débit volumique à partir de la différence de niveaux mesurée dans le manomètre.

\end{enumerate}
\end{document}